% Chapter 7, Exercises 11--15.

\begin{exercise}[11]
  \label{ex:ch7-11}
  Let $K$ be a global field. For any idele class character
  $\chi = (\chi_v)$, put
  \[
    W(\chi) = \prod_v W(\chi_v).
  \]
  \begin{enumerate}[label=(\alph*)]
    \item Show that $W(\chi)W(\chi^{-1}) = 1$. In particular, if $\chi$ is
      unitary, then $\abs{W(\chi)} = 1$.
    \item Let $K$ be a number field, and let $\chi$ be quadratic or trivial.
      Conclude from part~(a) that $W(\chi) = \pm 1$. [This number is called
      the \emph{sign of the functional equation} of $L(s, \chi)$.] Show,
      moreover, that if $\chi$ is unramified everywhere, then
      \[
        W(\chi) = \widetilde{\chi}(\calD_K)
      \]
      where $\widetilde{\chi}$ is the associated character of the class group
      $\Cl_K$ (cf.\ \RV{Proposition 5-19}). [\emph{Hint}: Use the previous
      problem.]
    \item Let $\chi$ be a quadratic idele class character, and let $E/K$ be
      the quadratic extension corresponding to the open subgroup
      $\ker(\chi)$ of $C_K$ by class field theory (\RV{Theorem 6-6}). Note
      that $\zeta_E(s) = \zeta_K(s) L(s, \chi)$ and deduce the formula
      \[
        \varepsilon(s, \chi) = \bigl(\abs{d_K}/\abs{d_E}^{1/2}\bigr)^{1-2s}
      \]
      where $d_K$ and $d_E$ are the discriminants of $K$ and $E$,
      respectively.
    \item (Hecke's Theorem, Serre's Proof) Let $K$ be a number field. Then
      prove:
      \begin{quote}
        \textsc{Theorem} (Hecke). \emph{The ideal class of $\calD_K$ is a
        square in the class group $\Cl_K$.}
      \end{quote}
      [\emph{Hint}: Observe that it suffices to show that
      $\widetilde{\chi}(\calD_K) = 1$ for every quadratic character
      $\widetilde{\chi}$ of $\Cl_K$. Then appeal to parts~(b) and~(c) of this
      problem.]
  \end{enumerate}
\end{exercise}

\begin{exercise}[12]
  \label{ex:ch7-12}
  We consider here the Fourier transforms of Schwartz--Bruhat functions.
  \begin{enumerate}[label=(\alph*)]
    \item Let $F$ be a local field. Show that for every $f \in \calS(F)$, its
      Fourier transform likewise lies in $\calS(F)$. [\emph{Hint}: For $F$
      Archimedean, this is a well-known classical fact. In the
      non-Archimedean case, use that $f$ is a linear combination of
      characteristic functions of the basic compact sets $P^n$, where $P$ is
      the unique prime associated with $F$.]
    \item Let $K$ be a global field and assume that $f \in \calS(\AAA_K)$.
      Show that the Fourier transform of $f$ is likewise in $\calS(\AAA_K)$.
      [\emph{Hint}: First prove that $\calS(\AAA_K)$ is generated by
      factorizable functions $\bigotimes_v f_v$, where each $f_v$ lies in
      $\calS(K_v)$ for all $v$, and then use the previous part to establish
      the result in this special case.]
    \item Use the results of \RV{Chapter 3} to show that the Fourier
      transform map $\calS(\AAA_K) \to \calS(\AAA_K)$ extends to an isometry
      $L^2(\AAA_K) \to L^2(\AAA_K)$.
  \end{enumerate}
\end{exercise}

\begin{exercise}[13]
  \label{ex:ch7-13}
  (Hecke Characters) Let $K$ be a number field of degree $d$, let $S$ be a
  finite set of places containing $S_\infty$, and let $J_K(S)$ be the group
  of fractional ideals prime to $S$. Furthermore, let $\II_K(S)$ denote the
  subgroup of $\II_K$ consisting of ideles $y = (y_v)$ such that $y_v = 1$ at
  every place $v$ in $S$.
  \begin{enumerate}[label=(\alph*)]
    \item Define $\alpha \colon \II_K(S) \to J_K(S)$ by
      \[
        \alpha(y) = \prod_{v \notin S} P_v^{v(y_v)}.
      \]
      Show that $\alpha$ is a surjective homomorphism with kernel
      \[
        U(S) = \{y \in \II_K(S) : y_v \in \gotho_v^\times \
        \forall v \text{ finite}\}.
      \]
      Conclude that every character $\chi$ of $\II_K$ that is unramified
      outside $S$ defines a character $\widetilde{\chi}$ of $J_K(S)$.
    \item Let $\beta$ be a \emph{Gr\"ossencharakter} (or \emph{Hecke
      character}) of $K$, which is to say a homomorphism from $J_K(S)$ to
      $\CC^\times$ for which there exists an integral ideal $M$ with support
      $S$, complex numbers $s_1, \ldots, s_d$ and integers
      $m_1, \ldots, m_d$ such that for every $\alpha \in K^\times(M)$, one
      has
      \[
        \beta((\alpha))
        = \prod_{j=1}^{d} \sigma_j(\alpha)^{m_j} \abs{\sigma_j(\alpha)}^{s_j}
      \]
      where $\{\sigma_j\}$ is the set of embeddings of $K$ in $\CC$. Show
      that every $\beta$ is of the form $\widetilde{\chi}$ for some character
      $\chi \colon \II_K \to \CC^\times$ that is trivial on $K^\times$.
      [\emph{Hint}: Show that $\beta$ defines a character of the ray class
      group $\Cl_K(M)$ and then lift it to the idele class group $C_K$.]
    \item Let $\Phi = \widetilde{\chi}$ for an idele class character $\chi$
      unramified outside $S$. Show that $\Phi$ is in fact a
      Gr\"ossencharakter. [\emph{Hint}: Take
      \[
        M = \prod_{v \in S} P_v^{\max(1, n_v)}
      \]
      where $n_v$ is the exponent of the conductor of $\chi_v$. For
      $\alpha \in K^\times(M)$, show that
      $\Phi((\alpha)) = \prod_{v \notin S} \chi_v(\alpha)
      = 1/(\chi_\infty \circ \sigma)(\alpha)$
      by using first that $\abs{\alpha}_v = 1$ for all $v$ in $S$ and second
      that $\chi$ is trivial on $K^\times$.]
  \end{enumerate}
\end{exercise}

\begin{exercise}[14]
  \label{ex:ch7-14}
  Fix an integer $k \ge 1$, and consider the \emph{polylogarithm} function
  \[
    l_k(z) = \sum_{n \ge 1} \frac{z^n}{n^k}.
  \]
  Note that $l_1(z) = -\log(1-z)$. The special case of $l_2(z)$ is called
  \emph{Euler's dilogarithm} function.
  \begin{enumerate}[label=(\alph*)]
    \item Show that the series is normally convergent in the open unit disk
      $\{\abs{z} < 1\}$ and that it has finite limit on the unit circle
      $S^1$ (respectively, $S^1 \smallsetminus \{1\}$) if $k > 1$
      (respectively, $k = 1$).
    \item Let $\chi$ be a Dirichlet character of conductor $m \ge 1$. For $w$
      an $m$th root of unity, put
      \[
        G(\chi, w) = \sum_{c \bmod m} \chi(c) w^c.
      \]
      Let $L(s, \chi)$ be the Dirichlet $L$-series, and assume that $\chi$ is
      not the trivial character if $k = 1$. Put $\chi(-1) = (-1)^r$. Then
      show that
      \[
        L(k, \chi) = \frac{1}{m} \sum_{1 \le b \le m/2}
        G(\chi, e^{-2\pi i b/m})
        \bigl[l_k(e^{2\pi i b/m}) + (-1)^r l_k(e^{-2\pi i b/m})\bigr].
      \]
      [\emph{Hint}: Expand $\chi$ in terms of the basis
      $\{e^{2\pi i b/m} : 0 \le b < m\}$ of $\CC[\ZZ/m\ZZ]$.]
    \item Show that
      \[
        G(\chi, e^{-2\pi i b/m}) = \overline{\chi}(b) G(\chi, e^{-2\pi i/m}).
      \]
      Simplify part~(b) accordingly for $k = 1$ and deduce
      \RV{Proposition 7-26}.
    \item For any $k \ge 1$, show that there is an elementary expression for
      the Dirichlet series $L(k, \chi)$ if $\chi(-1) = (-1)^k$.
  \end{enumerate}
\end{exercise}

\begin{exercise}[15]
  \label{ex:ch7-15}
  Let $\chi$ be a Dirichlet character of conductor $m \ge 1$. As above, put
  \[
    W(\chi) = \varepsilon(\chi\abs{\cdot}^{1/2})
    \quad \bigl(= \varepsilon(\tfrac12, \chi)\bigr).
  \]
  Show that
  \[
    W(\chi) = \frac{G(\chi)}{\sqrt{m}} (-i)^r
  \]
  where $r = 0$ (respectively, $1$) if $\chi(-1) = 1$ (respectively, $-1$),
  and
  \[
    G(\chi) = \sum_{a \bmod m} \chi(a) e^{2\pi i a/m}.
  \]
\end{exercise}
